\chapterbackmatter{Supplementary Exercises}

\begin{enumerate}
\SupplementaryExercise{1}{One Phase-Space Time Slice and Integrating Out a Quadratic Momentum}
  {(Adapted from Hong Liu and MIT OpenCourseWare, 8.323 Assignment 5, Problem 1.)}
  {mit-823-s23-ps5-ch9-curated-parent-1}
\begin{enumerate}
\item
For \(H=P^2/(2m)+V(Q)\), insert a momentum resolution of the identity
between two nearby position eigenstates.  Derive the short-time
phase-space kernel through first order in the time step, including its
normalization and the endpoint at which \(V\) is evaluated.
\item
Starting with the time-sliced phase-space path integral for
\(H=P^2/(2m)+V(Q)\), perform every momentum integral.  Derive the
configuration-space action and explain why the resulting determinant is
field independent in this example.
\end{enumerate}

\SupplementaryExercise{2}{The Schrödinger Equation from One Step}
  {(Adapted from Hong Liu and MIT OpenCourseWare, 8.323 Assignment 5, Problem 2.)}
  {mit-823-s23-ps5-ch9-curated-parent-2}
Propagate a smooth wavefunction through one infinitesimal
configuration-space time slice.  Expand consistently in the short time and
prove that the path-integral kernel generates the Schrödinger equation with
the correct kinetic and potential terms.  Track the Gaussian moments that
survive through first order and state how the endpoint prescription for the
potential enters the result.

\SupplementaryExercise{3}{The Harmonic-Oscillator Propagator from the Path Integral}
  {(Adapted from Daniel Harlow, \emph{Lecture Notes for Physics 8.323: Relativistic Quantum Field Theory I} (2024), Section 5.8, Homework Problem 3, printed p.~68.)}
  {harlow-qft1-2024-s5p8-problem3-v1ch09-s03}
Use the path integral to calculate the propagator
\[
 K_\omega(q',T;q,0)=\braket{q'}{e^{-\ii HT}q}
\]
of the simple harmonic oscillator
\[
 H=\frac{P^2}{2m}+\frac{k}{2}Q^2,
 \qquad \omega=\sqrt{\frac{k}{m}},
 \qquad T>0.
\]
The endpoints are fixed: \(q(0)=q\) and \(q(T)=q'\).  For comparison, the
preceding free-particle problem gives
\[
 K_0(q',T;q,0)
 =\sqrt{\frac{m}{2\pi \ii T}}
 \exp\left[\frac{\ii m(q'-q)^2}{2T}\right],
\]
where the square-root phase is selected by the short-time
\(i0\) prescription.

Write each path as a classical solution plus a fluctuation,
\(q(t)=q_{\rm cl}(t)+\eta(t)\), with
\(\eta(0)=\eta(T)=0\).  Expand \(\eta\) in its complete Fourier sine
series and perform the Gaussian integrals over all mode coefficients.
You may first keep only factors that depend on \(k\); fix every
\(k\)-independent factor at the end by requiring the result to approach
the displayed free-particle kernel as \(k\to0\).  Determine both the full
classical exponential and the determinant prefactor.

\SupplementaryExercise{4}{The Cubic Anharmonic Oscillator through Order \(\lambda^2\)}
  {(Adapted from University of Cambridge, Mathematical Tripos Part III,
  Advanced Quantum Field Theory, Paper 304 (2023), Question 1, p.~2,
  parts (a)--(c).)}
  {cambridge-partiii-2023-aqft-paper304-q1}
Use the source convention
\[
 Z[J]=\int\mathcal D x\,
 \exp\left[-\ii S[x]-\ii\int \dd t\,J(t)x(t)\right],
 \qquad
 x(t)\longleftrightarrow \ii\frac{\delta}{\delta J(t)}.
\]
The action for a harmonic oscillator of unit mass is
\[
 S[x]=\int \dd t\left(\frac12\dot x^{\,2}(t)
 -\frac12\omega^2x^2(t)\right),
 \qquad \omega>0.
\]
\begin{enumerate}
\item
Show, by Fourier transform or otherwise, that the propagator is
\[
 D(t-t')=\frac{1}{2\omega}e^{\ii\omega|t-t'|}.
\]
By introducing the source \(J(t)\), show that the harmonic-oscillator
generating functional can be written
\[
 Z_0[J]=Z_0[0]\exp\left[-\frac12\int \dd t\int \dd t'\,
 J(t)D(t-t')J(t')\right].
\]
\item
Now add the anharmonic term
\[
 S_\lambda[x]=-\frac{\lambda}{3!}\int \dd t\,x^3(t)
\]
to the action.  Write the interacting generating functional \(Z[J]\) as an
infinite series involving \(Z_0[J]\) and its functional derivatives.  Briefly
comment on the convergence of this series.
\item
Write down the Feynman rules for the anharmonic oscillator.  Hence find an
integral expression for the connected correlation function
\[
 \bra{\Omega}T\{x(t_1)x(t_2)\}\ket{\Omega}_{\rm conn}
\]
through order \(\lambda^2\), including every connected contribution at that
order.  Without evaluating the integrals, do you expect this correlation
function to diverge?  Briefly explain your reasoning, including its
ultraviolet behavior.
\end{enumerate}

\SupplementaryExercise{5}{Vacuum Survival under an External Force}
  {(Adapted from Mark Srednicki, Quantum Field Theory, Chapter 7, Problem 7.4.)}
  {srednicki-qft-path-integrals-ch9-curated-parent-2}
A ground-state oscillator is acted on by a real force \(f(t)\) of compact
support.  Evaluate the Gaussian path integral and express its vacuum
survival probability entirely in terms of the on-shell Fourier component
of the force.  Separate the phase from the modulus, identify the mean number
of produced quanta, and verify the adiabatic no-production limit.

\SupplementaryExercise{6}{The Free Scalar Source Functional and Wick Pairings by Differentiation}
  {(Inspired by Mark Srednicki, Quantum Field Theory, Chapter 8, Sections 8.1--8.2 | Mark Srednicki, Quantum Field Theory, Chapter 8, Eqs. 8.14--8.17.)}
  {srednicki-qft-path-integrals-ch9-curated-parent-3}
\begin{enumerate}
\item
For a real massive scalar coupled linearly to a source \(J\), complete the
functional square and derive \(Z_0[J]/Z_0[0]\).  State the momentum-space
kernel and its \(\ii\epsilon\) prescription in the mostly-plus convention.
\item
Differentiate the free scalar source functional four times and then set
the source to zero.  Derive the three pairings of the four-point function
and generalize the counting to \(2n\) insertions.
\end{enumerate}

\SupplementaryExercise{7}{The Propagator Is a Green Function}
  {(Adapted from John McGreevy, Physics 215A QFT, Fall 2023, Assignment 5,
  Problem 2, ``The propagator is a Green's function,'' pp.~1--2,
  parts (a)--(d).)}
  {mcgreevy-215a-f23-hw05-p2}
Work in \(d+1\) spacetime dimensions with the volume convention
\[
 \Box_x\equiv\partial_{x^\mu}\partial_x^{\mu}
 =\nabla_{\mathbf x}^2-\partial_{x^0}^2.
\]
For reference, the free real scalar has the mode expansion
\[
 \phi_0(x)=\int\frac{\dd^d\mathbf p}{(2\pi)^d\sqrt{2\omega_{\mathbf p}}}
 \left[a_{\mathbf p}e^{\ii p\cdot x}
 +a_{\mathbf p}^\dagger e^{-\ii p\cdot x}\right],
 \qquad \omega_{\mathbf p}=\sqrt{\mathbf p^2+m^2},\quad p^0=\omega_{\mathbf p},
\]
with
\[
 [a_{\mathbf p},a_{\mathbf q}^\dagger]
 =(2\pi)^d\delta^d(\mathbf p-\mathbf q),
 \qquad a_{\mathbf p}\ket{\Omega}=0.
\]
\begin{enumerate}
\item
Consider the retarded propagator
\[
 D_R(x-y)\equiv\theta(x^0-y^0)
 \bra{\Omega}[\phi(x),\phi(y)]\ket{\Omega}.
\]
Show that it is a Green function for the Klein--Gordon operator in the
sense that
\[
 (\Box_x-m^2)D_R(x-y)=a\,\delta^{d+1}(x-y),
\]
and find \(a\).
\item
Add the external-source term
\[
 \mathcal L= -\frac12\partial_\mu\phi\,\partial^\mu\phi
 -\frac12m^2\phi^2+\phi(x)j(x),
\]
where \(j(x)\) is a fixed \(c\)-number function.  Use the preceding result to
generalize the free mode expansion to a solution of the sourced field
equation with retarded boundary conditions.
\item
The time-ordered propagator is
\[
\begin{split}
 D(x-y)&\equiv
 \bra{\Omega}T\{\phi(x)\phi(y)\}\ket{\Omega}\\
 &=\theta(x^0-y^0)\bra{\Omega}\phi(x)\phi(y)\ket{\Omega}
 +\theta(y^0-x^0)\bra{\Omega}\phi(y)\phi(x)\ket{\Omega}.
\end{split}
\]
Show that it too is a Green function for the Klein--Gordon operator, and
determine its delta-function normalization.

\emph{Hint.}  Use the canonical equal-time commutation relations
\[
 [\phi(x^0,\mathbf x),\phi(x^0,\mathbf y)]=0,
 \qquad
 [\partial_{x^0}\phi(x^0,\mathbf x),\phi(x^0,\mathbf y)]
 =-\ii\delta^d(\mathbf x-\mathbf y).
\]
Do not neglect \(\partial_t\theta(t)=\delta(t)\): the time derivatives act
on the time-ordering symbol.
\item
This is an example of a fact that is easier to understand from the path
integral: explain why correlation functions solve the field equations only
up to contact terms, and relate that statement to the delta function found
in the preceding part.
\end{enumerate}

\SupplementaryExercise{8}{The Scalar Vacuum Wave Functional}
  {(Adapted from Mark Srednicki, Quantum Field Theory, Chapter 8, Problem 8.8.)}
  {srednicki-qft-path-integrals-ch9-curated-parent-5}
Treat each spatial Fourier mode of a free real scalar as an oscillator.
Derive the vacuum wave functional, its nonlocal position-space kernel, and
the equal-time covariance of two field modes.  Fix the normalization as a
regulated product of oscillator normalizations and identify which part
survives in normalized expectation values.

\SupplementaryExercise{9}{Schwinger--Dyson Equations from Functional Integration by Parts}
  {(Adapted from J. McGreevy, Physics 215A Quantum Field Theory, Fall 2023, Assignment 5, Problem 3, printed pp.~2--3.)}
  {mcgreevy-215a-f23-hw05-p3-v1ch09-s09}
Consider the Lorentzian path integral
\[
 Z=\int\mathcal D\phi\,e^{\ii S[\phi]},
\]
defined with the usual \(i0\) prescription.  Invariance of the measure
under a change of field variable implies
\[
 0=\int\mathcal D\phi\,
 \frac{\delta}{\delta\phi(x)}\bigl(\hbox{anything}\bigr),
\]
provided the differentiated expression has no boundary contribution at
large field.  To explain this rather than merely assume it, discretize
spacetime.  Then
\[
 \int\mathcal D\phi=\int\prod_zd\phi_z,
 \qquad \phi_z\equiv\phi(z),
\]
and the identity is just ordinary integration by parts in the variable
\(\phi_x\):
\[
 \int \dd\phi_x\,\frac{\partial}{\partial\phi_x}
 \bigl(\hbox{anything}\bigr)=0.
\]
These functional integration-by-parts identities are the
Schwinger--Dyson equations.  They are nonperturbative and implement the
field equations inside correlation functions.
\begin{enumerate}
\item
Evaluate
\[
 0=\int\mathcal D\phi\,
 \frac{\delta}{\delta\phi(x)}
 \left(\phi(y)e^{\ii S[\phi]}\right)
\]
and prove the normalized identity
\[
 \boxed{
 \frac1Z\int\mathcal D\phi\,e^{\ii S[\phi]}
 \frac{\delta S}{\delta\phi(x)}\phi(y)
 =\ii\delta^{d+1}(x-y).}
\]
\item
Apply this result to the free massive real scalar action in the volume
convention,
\[
 S_0[\phi]=\int d^{d+1}x\left(
 -\frac12\partial_\mu\phi\partial^\mu\phi
 -\frac12m^2\phi^2\right).
\]
Show that the time-ordered two-point function obeys the classical field
equation away from the other insertion and derive its contact term:
\[
 \boxed{
 (\Box_x-m^2)
 \bra{\Omega}T\{\phi(x)\phi(y)\}\ket{\Omega}
 =\ii\delta^{d+1}(x-y).}
\]
\item
Derive the corresponding Schwinger--Dyson equation for the time-ordered
three-point function of the free field.  Display every contact term and
then simplify the result using the free-vacuum one-point function.
\end{enumerate}

\SupplementaryExercise{10}{Low-Order Correlators in \(\lambda\phi^4\)}
  {(Adapted from Hong Liu and MIT OpenCourseWare, 8.323 Assignment 6, Problem 2.)}
  {mit-823-s23-ps6-ch9-curated-parent-2}
\begin{enumerate}
\item Enumerate every Wick topology contributing to the normalized
two-point function through order \(\lambda^2\).  Give its position-space
integral and symmetry factor, showing the cancellation of vacuum factors.
\item Do the same for the connected four-point function through order
\(\lambda^2\), including the three one-loop channels and the one-vertex
tree term.
\end{enumerate}

\SupplementaryExercise{11}{Vacuum Energy and Connected Vacuum Graphs}
  {(Adapted from Hong Liu and MIT OpenCourseWare, 8.323 Assignment 6, Problem 3.)}
  {mit-823-s23-ps6-ch9-curated-parent-3}
\begin{enumerate}
\item Show that translational invariance makes
\(W_0=-\ii\log Z[0]\) proportional to spacetime volume and identify the
coefficient with the vacuum-energy density.
\item Prove combinatorially that \(W_0\) is the sum of connected vacuum
graphs, with one inverse symmetry factor for each topology.
\item In \(\lambda\phi^4/4!\) theory, calculate the vacuum-energy-density
shift through order \(\lambda^2\), displaying all connected integrals and
their symmetry factors.
\end{enumerate}

\SupplementaryExercise{12}{Cancellation of Arbitrary Vacuum Bubbles}
  {(Adapted from Hong Liu and MIT OpenCourseWare, 8.323 Assignment 6, Problem 4.)}
  {mit-823-s23-ps6-ch9-curated-parent-4}
For a general polynomial scalar interaction, factor diagrams contributing
to an unnormalized \(n\)-point function into vacuum and source-connected
components.  Prove to all perturbative orders that normalization by
\(Z[0]\) cancels every vacuum component.  State the argument as a generating
functional identity and then interpret the same cancellation graph by
graph, including symmetry factors.

\SupplementaryExercise{13}{A Wave Functional as a Euclidean Boundary and A Parity-Twisted Euclidean Trace}
  {(Inspired by David Skinner, Advanced Quantum Field Theory, Sections 1.2.3 and 3.1 | David Skinner, Advanced Quantum Field Theory, Example Sheet 1, Problem 6.)}
  {skinner-aqft-path-integrals-2018-ch9-curated-parent-1}
\begin{enumerate}
\item
Represent \(\braket{\phi_f}{e^{-TH}\psi_i}\) as a Euclidean path integral
with one fixed boundary and one boundary weighted by an arbitrary initial
wave functional.  Show how the ground-state functional is projected out as
\(T\to\infty\).
\item
Find the Euclidean path-integral boundary condition that computes
\(\operatorname{Tr}(Pe^{-TH})\) for a particle on the line, where \(P\)
is parity.  Evaluate the trace for the harmonic oscillator and compare it
with the energy-basis answer.
\end{enumerate}

\SupplementaryExercise{14}{Pfaffian Covariance and Canonical Form}
  {(Adapted from David Skinner, Advanced Quantum Field Theory, Example Sheet 1, Problem 3.)}
  {skinner-aqft-path-integrals-2018-ch9-curated-parent-2}
For \(2N\) real Grassmann variables and an invertible antisymmetric matrix
\(A\), derive the transformation of the Berezin measure under
\(\theta=N\theta'\).  Use it to prove
\(\operatorname{Pf}(N^{\mathsf T}AN)=\det(N)\operatorname{Pf}(A)\).
Then reduce \(A\) by a change of basis to \(2\times2\) antisymmetric
blocks and use that canonical form to prove
\(\operatorname{Pf}(A)=\pm\sqrt{\det A}\).  Fix the sign by continuity
and check every formula for one Grassmann pair.

\SupplementaryExercise{15}{An Exact Four-Grassmann Interaction}
  {(Adapted from David Skinner, Advanced Quantum Field Theory, Example Sheet 1, Problem 4.)}
  {skinner-aqft-path-integrals-2018-ch9-curated-parent-3}
Evaluate exactly the partition function of four real Grassmann variables
with an antisymmetric quadratic kernel and a completely antisymmetric
quartic coupling.  Reproduce the same answer by enumerating the allowed
vacuum contractions.  Track the ordering convention for the Berezin measure,
express the quadratic contribution as a Pfaffian, and check both the free and
pure-quartic limits.

\SupplementaryExercise{16}{Bosonic and Fermionic Determinant Powers and Fermionic Wick's Theorem from a Pfaffian}
  {(Inspired by David Skinner, Advanced Quantum Field Theory, Sections 2.2 and 2.4 | David Skinner, Advanced Quantum Field Theory, Sections 2.4 and 3.2.)}
  {skinner-aqft-path-integrals-2018-ch9-curated-parent-4}
\begin{enumerate}
\item
Compare a finite-dimensional real bosonic Gaussian, a complex Grassmann
Gaussian, and a real Grassmann Gaussian.  Derive respectively the
determinant powers \(-1/2\), \(+1\), and the Pfaffian, including the
assumptions on their kernels.
\item
Introduce Grassmann sources for a real fermionic Gaussian.  Complete the
square, differentiate four times, and derive the signed sum over the three
pairings.  Explain how the result generalizes to \(2N\) insertions.
\end{enumerate}

\SupplementaryExercise{17}{Berezin Gaussians, Fermionic Sources, and Determinant Signs}
  {(Adapted from Hong Liu and MIT OpenCourseWare, 8.323 Assignment 10, Problem 3.)}
  {mit-823-s23-ps10-grassmann-ch9-curated-parent-1}
\begin{enumerate}
\item
With a stated order for both variables and differentials, evaluate
\(\int \dd\bar\theta\,\dd\theta\,e^{-\bar\theta a\theta}\) directly by
expansion.  Check translation invariance under independent Grassmann
shifts.
\item
For \(N\) independent pairs \(\theta_i,\bar\theta_i\), diagonalize or
triangularize the kernel and prove that their Gaussian integral equals
\(\det A\).  Explain why the determinant is not inverted as it is for
commuting variables.
\item
Insert \(\theta_k\bar\theta_l\) into the complex Grassmann Gaussian.
Derive the result by differentiating the determinant and identify carefully
the ordering of the indices of \(A^{-1}\).
\item
Evaluate the normalized Gaussian expectation of
\(\theta_i\bar\theta_j\theta_k\bar\theta_l\).  Derive the relative minus
sign between its two pairings directly from Grassmann source
differentiation.
\item
Add independent Grassmann sources to the free Dirac action.  Shift
\(\psi,\bar\psi\), derive the normalized generating functional, and recover
the momentum-space Dirac propagator in the volume's gamma-matrix
convention.
\item
Integrate out a Dirac field coupled bilinearly to a background matrix
\(\Gamma(x)\).  Expand the resulting determinant and show explicitly why
each connected closed fermion loop carries one additional minus sign.
\end{enumerate}

\SupplementaryExercise{18}{The Complete Gaussian Generating Function}
  {(Adapted from John McGreevy, Physics 215A, Fall 2023, Assignment 4, Problem 2.)}
  {mcgreevy-215a-f23-hw04-p2}
Let \(A\) be a real symmetric positive-definite \(N\times N\) matrix.
\begin{enumerate}
\item
Derive the one-dimensional source integral
\[
 Z_a(j)=\int_{-\infty}^{\infty}\dd x\,
 e^{-ax^2/2+jx}
\]
including its normalization.  Establish the source-free Gaussian integral
without assuming its value, and state the convergence condition.
\item
Diagonalize \(A\) and evaluate
\[
 Z_A(J)=\int_{\mathbb R^N}d^Nx\,
 e^{-x^{\mathsf T}Ax/2+J^{\mathsf T}x}
\]
in terms of \(\det A\) and \(A^{-1}\).  Explain why an orthogonal change of
basis introduces no additional Jacobian magnitude.
\item
With normalized expectation values defined by \(Z_A(0)\), obtain
\(\langle x_i x_j\rangle\) both by differentiating \(Z_A\) and by
differentiating \(\log Z_A\).  Identify what derivatives of the logarithm
generate, and state the continuum functional version of the result.
\end{enumerate}

\SupplementaryExercise{19}{The Characteristic Function of a Gaussian Quantum Coordinate}
  {(Adapted from John McGreevy, Physics 215A, Fall 2023, Assignment 4, Problem 3.)}
  {mcgreevy-215a-f23-hw04-p3}
Let \(q\) be a coordinate in the vacuum of a centered Gaussian quantum
system, so its Hamiltonian is quadratic and
\(\bra{\Omega}q\ket{\Omega}=0\).  Determine the full \(K\)-dependence of
\(\bra{\Omega}e^{\ii Kq}\ket{\Omega}\) in terms of
\(\bra{\Omega}q^2\ket{\Omega}\).  Derive the identity once from the
Gaussian position-space measure and once from oscillator operator algebra.
Then give the corrected formula when the Gaussian has nonzero mean
\(\bar q\).
\end{enumerate}
